\section{Related Work}

\textbf{DLMs.}
DLMs generate text by iteratively denoising discrete token sequences. Prior work explored discrete-state diffusion, continuous diffusion in embedding space, and diffusion with pretrained masked language models~\citep{austin2021structured,li2022diffusionlm,he2023diffusionbert}. Masked discrete diffusion provided a scalable formulation, later demonstrated at scale by LLaDA and Dream~\citep{sahoo2024simple,nie2025neurips-large,ye2025dream}. Fully parallel masked DLMs predict multiple positions per step, enabling bidirectional context modeling and parallel generation. However, repeatedly processing the full sequence incurs redundant computation on largely unchanged tokens, while bidirectional attention prevents direct reuse of the append-only KV cache used by autoregressive models~\citep{wu2025fastdllm,liu2025dllmcache}.

\textbf{From DLMs to BDLMs.}
Fast-dLLM demonstrates the potential of block-level organization for DLM acceleration through block-wise approximate KV caching and confidence-aware parallel decoding~\citep{wu2025fastdllm}. Subsequent methods explore adaptive cache reuse at different granularities, including segment-level updates, delayed updates, and token-level state selection~\citep{liu2025dllmcache,ma2025dkvcache,jiang2025d2cache}. BDLMs make this structure part of the generation paradigm: Block Diffusion generates blocks autoregressively while refining tokens through diffusion within each block~\citep{arriola2025iclr-block}; SDAR adapts pretrained autoregressive models to block-wise diffusion~\citep{cheng2025sdar}; and LLaDA2.0 and LLaDA2.1 extend the scale and token-editing capabilities of this model family~\citep{nie2025llada20,bie2026llada21speedingtextdiffusion}. Fast-dLLM v2 further combines block-diffusion model design with hierarchical caching and parallel decoding~\citep{wu2025fastdllmv2}.

\textbf{Inference optimization for existing BDLMs.}
Recent work optimizes existing BDLMs, including SDAR and LLaDA2.x, by reducing current-block computation or historical-prefix attention~\citep{liang2026focus,xi2026losa}. FOCUS estimates token-wise decodability to skip positions unlikely to be updated, but under confidence-based decoding this may increase the number of denoising steps, limiting end-to-end gains~\citep{liang2026focus}. LoSA sparsifies long-context prefix attention by recomputing attention for active queries and reusing outputs for stable ones, while repeated query-dependent KV selection still incurs nontrivial overhead~\citep{xi2026losa}. As a result, the reduction in attention computation does not always translate into substantial end-to-end speedup.